\chapter{Related Work}\label{chap:research}

To properly contextualize the operational environment of this thesis, it is necessary to examine the fundamental domain of continuous black-box optimization, the algorithmic vulnerabilities exposed by stochastic noise, and the recent shift toward Automated Algorithm Design. Continuous black-box optimization focuses on identifying the global minimum of a function $f: \mathbb{R}^D \rightarrow \mathbb{R}$ where the analytical gradient and any internal structural information about the function are unavailable, leaving sampling queries as the primary mechanism for obtaining information. In standard literature, algorithms designed for these problems are commonly evaluated against the deterministic Black-Box Optimization Benchmark (BBOB) suite \cite{Hansen2010}. The BBOB framework also standardized the use of empirical cumulative distribution functions (ECDFs) for characterizing algorithm performance, measuring runtime by the number of function evaluations required to achieve a specific target precision ($\Delta f_t$). Over recent decades, human-designed optimization methods have produced robust heuristics for navigating these deterministic spaces, including the Nelder-Mead simplex algorithm \cite{Nelder1965}, Differential Evolution (DE) \cite{Storn1997}, Particle Swarm Optimization (PSO) \cite{Kennedy1995}, and the Covariance Matrix Adaptation Evolution Strategy (CMA-ES) \cite{Hansen2001}. Empirical benchmarks on the noiseless BBOB testbed have consistently established CMA-ES as a strong baseline for continuous optimization, outperforming local search methods such as NEWUOA \cite{Powell2006} on highly conditioned, non-separable, and ill-conditioned multi-modal topologies due to its covariance adaptation mechanism \cite{Hansen2010}. However, these methods rely on information derived from observed search behavior and assumptions about the objective landscape, which can make their performance sensitive to changes in the evaluation environment.

While human-engineered algorithms can perform well on deterministic functions, practical industrial applications are often inherently stochastic. Objective evaluations in these environments can be affected by sensor inaccuracies, which are systematically modeled in continuous optimization testbeds, most notably the BBOB-noisy suite \cite{Hansen2009Noisy}. As established in the foundational taxonomy of uncertain environments by Jin and Branke \cite{Jin2005}, sources of uncertainty in optimization can be classified according to whether they affect the parameter space, time space, or fitness space. This thesis focuses on noise acting directly within the fitness space, where the true time-invariant objective function $f(x)$ is affected by an error component $\epsilon$, making the observed objective value a random variable centered around the underlying true fitness. In such settings, conventional gradient-free heuristics can be affected because their selection mechanisms rely on ranking candidates according to observed performance, and fitness noise can therefore distort these rankings. Jin and Branke's evaluation showed that traditional approaches, such as static sample averaging or increasing population sizes, can reduce the variance caused by noise, but at the cost of slower optimization and increased function evaluation budgets \cite{Jin2005}. Consequently, a sub-optimal candidate may temporarily obtain a better position in the population because of a favorable statistical anomaly, causing the algorithm to adjust its covariance matrices or search steps based on misleading information \cite{Jin2005, Hansen2009Noisy}. Traditional parameter tuning, which relies on selecting fixed hyperparameter configurations before runtime, does not fully resolve this trade-off because fixed parameter choices cannot adapt to changes in local noise across the search space. This limitation motivates the investigation of whether autonomous logic generators can discover structural alternatives to static computational approaches.

To reduce the reliance on manual algorithm design, techniques such as hyperparameter tuning, automated algorithm selection, and traditional Automated Algorithm Design (AAD) were developed long before the advent of deep learning. Early AAD frameworks relied on approaches such as Genetic Programming \cite{Burke2013} or modular spaces designed by humans \cite{vanRijn2016}, but their predefined components and search spaces could limit the resulting discoveries. More recently, Large Language Models have been incorporated into AAD, with their potential for mathematical discovery demonstrated by Romera-Paredes et al. \cite{Romera2023}. Their FunSearch framework addresses the tendency of LLMs to produce incorrect or unverifiable outputs by combining the model with a systematic, closed-loop evaluator. Instead of asking the LLM to generate raw numerical solutions, the framework constrains it to evolve the critical logical components of a program skeleton. The authors demonstrated that this approach could produce novel and verifiable algorithmic constructions. FunSearch discovered new large-sized caps for the cap-set problem that improved on existing lower bounds established by human mathematicians, while also generating compact and reusable bin-packing heuristics that outperformed traditional greedy algorithms on large-scale datasets \cite{Romera2023}. These results demonstrate the potential of foundation models to generate new algorithmic constructions and provide a basis for applying similar closed-loop approaches to real-valued optimization problems.

Following this work, subsequent studies have investigated different ways of incorporating LLMs into optimization pipelines. For example, the EvoLLM framework \cite{Lange2024} uses an LLM directly as an evolutionary numerical sampler to predict coordinate distributions. While this approach can perform reasonably well, it does not provide the same interpretability and reusability offered by explicit algorithmic code, and it also requires a relatively large number of prompt tokens. In its evaluation, EvoLLM achieved reasonable performance on a subset of eight continuous BBOB functions, but the experiments were limited to a single 5-dimensional instance of each problem, which limits the conclusions that can be drawn about its scalability and generalizability \cite{Lange2024}. In contrast, frameworks such as Evolution of Heuristics (EoH) \cite{Liu2024} and Algorithm Evolution using LLMs (AEL) \cite{AEL2024} focus the LLM's generative role on producing reusable code blocks. These approaches achieved better performance and lower optimality gaps than traditional heuristics and deep reinforcement learning models, although their evaluations were mainly limited to small, discrete combinatorial problems such as the Traveling Salesperson Problem (TSP) \cite{AEL2024, Liu2024}. Related work on multi-agent language model structures, particularly the Reflexion framework \cite{Shinn2023} and automated self-debugging approaches \cite{Chen2023}, has also shown the value of providing execution traces and code tracebacks as feedback to the generating model. Chen et al.\ demonstrated that asking a model to perform zero-human-feedback "rubber duck debugging" by producing line-by-line natural language explanations alongside unit-test checks can improve functional accuracy by up to 12\% on programmatic translation tasks and improve success on complex logical problems by 9\%, while achieving performance comparable to baselines using a $10\times$ larger sampling budget \cite{Chen2023}. These results support the use of automated, text-based runtime feedback as a mechanism for guiding the generation and refinement of algorithmic code.

Building on this work on executable code generation and automated self-correction, this thesis extends the Large Language Model Evolutionary Algorithm (LLaMEA) framework \cite{LLaMEA2024}. LLaMEA was specifically developed to automate the discovery of continuous black-box optimization heuristics and implements the closed-loop concept by placing an LLM within a $1+1$ Evolution Strategy. The framework iteratively prompts the model to generate complete, reusable Python classes, executes the generated code on target mathematical problems using the IOHexperimenter benchmarking suite, and feeds runtime error tracebacks and aggregated performance metrics, specifically the Area Over the Convergence Curve (AOCC), back to the LLM as feedback for further refinement and debugging. Foundational experiments showed that LLaMEA-generated algorithms could discover novel heuristic logic that outperformed state-of-the-art metaheuristics, including vanilla CMA-ES and Differential Evolution, on the 5-dimensional continuous BBOB suite. The generated algorithms also maintained competitive generalizability when applied zero-shot to unseen 10-dimensional and 20-dimensional instances \cite{LLaMEA2024}. However, existing applications of the framework have focused on clean, noiseless functions, leaving its behavior under persistent statistical noise largely unexplored. If an algorithmic code class evolved under deterministic conditions is exposed to persistent fitness noise, its ability to track selection can be affected by corrupted rankings. By systematically injecting a controlled statistical noise profile into the evaluation sandbox, this thesis therefore extends the LLaMEA framework to investigate whether closed-loop evolutionary feedback can enable a language model to discover specialized, noise-resilient optimization heuristics.
